\documentclass[../../main.tex]{subfiles}
\begin{document}

{\bf[解]}

\begin{figure}[htb]
  \centering
  % 俯角 70°，方位角 120° で視認性を確保
  \tdplotsetmaincoords{70}{120}
  \begin{tikzpicture}[tdplot_main_coords, scale=3.2, >=stealth, every node/.style={font=\small}]

    % --- 1. 頂点座標の定義 ---
    % P(1, 0.5, 0), Q(1, -0.5, 0), R(0.25, 0, sqrt(3)/4 ≈ 0.4330)
    \coordinate (O) at (0, 0, 0);
    \coordinate (P) at (1.0,  0.5, 0.0);
    \coordinate (Q) at (1.0, -0.5, 0.0);
    \coordinate (R) at (0.25, 0.0, 0.4330);

    % --- 2. 座標軸 ---
    \draw[->, thick] (0, 0, 0) -- (1.4, 0, 0) node[below left=2pt] {$x$};
    \draw[->, thick] (0, -0.9, 0) -- (0, 0.9, 0) node[right=2pt] {$y$};
    \draw[->, very thick, red!80!black] (0, 0, -0.15) -- (0, 0, 0.8) node[above=2pt] {$z$ (回転軸)};
    \node[below left=2pt] at (O) {$O$};

    % 回転軸の回転矢印
    \tdplotdrawarc[->, thick, red!85!black]{(0,0,0.65)}{0.15}{30}{320}{anchor=south}{}

    % --- 3. xy 平面への投影・基準補助線 ---
    \draw[dotted, gray!60] (P) -- (1.0, 0, 0) -- (Q);
    \draw[dashed, gray!60] (R) -- (0.25, 0, 0);
    \draw[dotted, gray!60] (0.25, 0, 0) -- (1.0, 0, 0);
    \draw[dashed, gray!50] (R) -- (0, 0, 0.4330);
    \fill (0, 0, 0.4330) circle (0.4pt) node[left=1pt, font=\scriptsize] {$\frac{\sqrt{3}}{4}$};

    % --- 4. 正三角形の板 S (PQR) ---
    \fill[blue!15, opacity=0.75] (P) -- (Q) -- (R) -- cycle;
    \draw[very thick, blue!85!black] (P) -- (Q) -- (R) -- cycle;

    % 板の重心/内部ラベル
    \node[blue!90!black, font=\footnotesize] at ($ (P)!0.333!(Q)!0.333!(R) + (0.1, 0, 0.02) $) {$S$};

    % --- 5. 頂点プロットとラベル ---
    \fill (P) circle (0.7pt) node[right=2pt] {$P\left(1, \frac{1}{2}, 0\right)$};
    \fill (Q) circle (0.7pt) node[below left=1pt] {$Q\left(1, -\frac{1}{2}, 0\right)$};
    \fill (R) circle (0.7pt) node[above right=1pt] {$R\left(\frac{1}{4}, 0, \frac{\sqrt{3}}{4}\right)$};

  \end{tikzpicture}
  \caption{正三角形の板 $S$（頂点 $P, Q, R$）と回転軸（$z$ 軸）の空間配置}
  \label{fig:triangle_plate_3d}
\end{figure}

辺$PR$，$QR$の方程式はそれぞれ以下のようになる．
\begin{align}
  PR&:
  \begin{pmatrix}x\\y\\z
  \end{pmatrix}=
  \begin{pmatrix}1\\1/2\\0
  \end{pmatrix}+t
  \begin{pmatrix}-3/4\\-1/2\\\sqrt{3}/4
  \end{pmatrix} \\
  QR&:
  \begin{pmatrix}x\\y\\z
  \end{pmatrix}=
  \begin{pmatrix}1\\-1/2\\0
  \end{pmatrix}+t
  \begin{pmatrix}-3/4\\1/2\\\sqrt{3}/4
  \end{pmatrix}
\end{align}
これらと平面$z=k$の交点を$P_{k},Q_{k}$とすると，いずれも
\begin{align}
  k = \dfrac{\sqrt{3}t}{4} \quad \therefore t = \dfrac{4k}{\sqrt{3}}
\end{align}
に対応するから，
\begin{align}
  \begin{cases}
    p&=\dfrac{1}{2}\left(1-\dfrac{4\sqrt{3}}{3}k\right)  \\
    q&=1-\sqrt{3}k
  \end{cases}
\end{align}
として
\begin{align}
  P_{k}(q,p,k), Q_{k}(q,-p,k)
\end{align}
である．
故に，図形の$z=k(0\le k\le \dfrac{\sqrt{3}}{4})$での切断面は
\begin{align}
  &x=q \\
  &|y|\le p
\end{align}
となる．これを図示すると下図の太線部である．

\begin{figure}[htb]
  \centering
  \begin{tikzpicture}[scale=4, >=stealth, every node/.style={font=\small}]

    % --- 1. 定数・代表値 (k = 0.2 のとき q = 0.6536, p = 0.3614) ---
    \def\qVal{0.6536}
    \def\pVal{0.3614}

    \coordinate (O)      at (0, 0);
    \coordinate (Qmid)   at (\qVal, 0);
    \coordinate (Qtop)   at (\qVal, \pVal);
    \coordinate (Qbot)   at (\qVal, -\pVal);

    % --- 2. 座標軸 ---
    \draw[->, thick] (-0.2, 0) -- (1.1, 0) node[right] {$x$};
    \draw[->, thick] (0, -0.65) -- (0, 0.65) node[above] {$y$};
    \node[below left=2pt] at (O) {$O$};

    % --- 3. 座標補助線 (破線) ---
    \draw[dashed, gray!70] (0, \pVal) -- (Qtop);
    \draw[dashed, gray!70] (0, -\pVal) -- (Qbot);

    % --- 4. 切断面の線分 (x = q, |y| <= p) ---
    \draw[line width=2.4pt, blue!85!black] (Qbot) -- (Qtop);

    % --- 5. 点プロットとラベル ---
    \fill (Qmid) circle (0.6pt) node[below left=1pt] {$q$};
    \fill[red!85!black] (Qtop) circle (0.7pt) node[right=2pt, red!85!black] {$(q, p)$};
    \fill[red!85!black] (Qbot) circle (0.7pt) node[right=2pt, red!85!black] {$(q, -p)$};

    % y 軸上の目盛り
    \draw[gray!80!black] (-0.03, \pVal) -- (0.03, \pVal);
    \node[left=2pt, font=\scriptsize] at (0, \pVal) {$p$};
    \draw[gray!80!black] (-0.03, -\pVal) -- (0.03, -\pVal);
    \node[left=2pt, font=\scriptsize] at (0, -\pVal) {$-p$};

    % 領域の注釈
    \draw[<-, thin, blue!85!black] (\qVal, 0.18) to[bend left=18] (0.85, 0.35)
    node[right, font=\footnotesize, blue!90!black] {切断面: $x=q, \; |y| \le p$};

  \end{tikzpicture}
  \caption{平面 $z=k$ による板 $S$ の切断面（線分部）}
  \label{fig:cross_section_line}
\end{figure}

したがって，切断面上の点と原点との距離の最大小値はそれぞれ$(q,p)$および$(q,0)$で
\begin{align}
  r_{max}^2&=p^2+q^2 \\
  r_{min}^2&=q^2
\end{align}
である．故にこの平面での回転体の切断面は，下図のようにドーナツ型の領域になる．

\begin{figure}[htb]
  \centering
  \begin{tikzpicture}[scale=3.2, >=stealth, every node/.style={font=\small}]

    % --- 1. 定数・半径パラメータ (k = 0.2 の代表値) ---
    \def\qVal{0.6536}                             % x 座標 / 内径 r_min
    \def\pVal{0.3614}                             % y 半幅
    \def\rMax{0.7469}                             % 外径 r_max = sqrt(p^2 + q^2)

    \coordinate (O)     at (0, 0);
    \coordinate (Qmid)  at (\qVal, 0);
    \coordinate (Qtop)  at (\qVal, \pVal);
    \coordinate (Qbot)  at (\qVal, -\pVal);

    % --- 2. ドーナツ型（円環）領域の斜線ハッチング ---
    % 外円から内円をくり抜いた領域
    \fill[pattern=north east lines, pattern color=blue!55, opacity=0.7, even odd rule]
    (O) circle (\rMax)
    (O) circle (\qVal);

    % --- 3. 座標軸 ---
    \draw[->, thick] (-1.05, 0) -- (1.15, 0) node[right] {$x$};
    \draw[->, thick] (0, -1.05) -- (0, 1.15) node[above] {$y$};
    \node[below left=2pt] at (O) {$O$};

    % --- 4. 境界円 ---
    % 内円 (最小半径: 破線)
    \draw[thick, dashed, teal!85!black] (O) circle (\qVal);
    % 外円 (最大半径: 実線)
    \draw[very thick, blue!85!black] (O) circle (\rMax);

    % --- 5. 図2の元となる切断線分 (x = q, |y| <= p) の重畳表示 ---
    % 太い実線で円環内に配置
    \draw[line width=2.6pt, orange!90!black] (Qbot) -- (Qtop);

    % 線分上の重要点プロット
    \fill[teal!85!black] (Qmid) circle (0.7pt) node[below left=1pt, font=\scriptsize] {$(q, 0)$};
    \fill[red!85!black]  (Qtop) circle (0.7pt) node[above right=1pt, font=\scriptsize, red!85!black] {$(q, p)$};
    \fill[red!85!black]  (Qbot) circle (0.7pt) node[below right=1pt, font=\scriptsize, red!85!black] {$(q, -p)$};

    % --- 6. 半径ベクトルと距離の対応関係 ---
    % (i) 最小距離: 原点から (q, 0) へのベクトル (内円との接点)
    \draw[->, thick, teal!85!black] (O) -- (Qmid);
    \node[teal!95!black, font=\footnotesize, above=2pt] at ($ (O)!0.5!(Qmid) $) {$r_{\min} = q$};

    % (ii) 最大距離: 原点から端点 (q, p) への動径ベクトル (外円上の点)
    \draw[->, thick, red!85!black] (O) -- (Qtop);
    \node[red!90!black, font=\footnotesize, above left=1pt] at ($ (O)!0.55!(Qtop) $) {$r_{\max} = \sqrt{p^2+q^2}$};

    % 線分の引き出し線と注釈
    \draw[<-, thin, orange!90!black] ($ (Qtop)!0.5!(Qmid) $) to[bend left=20] (1.05, 0.28)
    node[right, font=\footnotesize, orange!90!black, align=left] {回転させる線分\\($x=q, \, |y| \le p$)};

  \end{tikzpicture}
  \caption{平面 $z=k$ での回転体の切断面（円環・斜線部）}
  \label{fig:cross_section_annulus_with_segment}
\end{figure}

この面積$S(k)$は
\begin{align}
  S(k)=\pi(r_{max}^2- r_{min}^2) =\pi p^2
\end{align}
となる．従って求める体積$V$は
\begin{align}
  V
  &=\int_0^{\sqrt{3}/4}S(k)dk \\
  &=\dfrac{\pi}{4}\int_0^{\sqrt{3}/4}\left(1-\dfrac{4\sqrt{3}}{3}k\right)^2dk \\
  &=\dfrac{\pi}{4}\left[\dfrac{\sqrt{3}}{12}\left(\dfrac{4\sqrt{3}}{3}k-1\right)^3\right]_0^{\sqrt{3}/4} \\
  &=\dfrac{\sqrt{3}\pi}{48}\cdots\text{(答)}
\end{align}
となる．
\newpage
\end{document}
